% GENERATED FILE. Edit canonical lessons, then run npm run generate:books.
% canonical-source-hash: fnv1a64:1aa6026f8500fe57
% canonical-lessons: ZH-C269-congxiao, ZH-C269-xian, ZH-W269-xian-first, ZH-C269-tongnian, ZH-C269-benlai, ZH-C269-wanju

\chapter{When I Was Little}
\label{ch:zh-when-i-was-little}
\hlchaptermodality{269}

\begin{chapteropening}
\textbf{By the end of this chapter:} \emph{I can say what I have done since childhood, what comes first, what my childhood was like, what I meant to do, and talk about toys.}

Complete the last lesson of chapter 269: I can say what I have done since childhood, what comes first, what my childhood was like, what I meant to do, and talk about toys.
\end{chapteropening}

\section[\texorpdfstring{\zh{从小} (cóngxiǎo)}{从小 (cóngxiǎo)}]{\zh{从小} (cóngxiǎo) — since childhood; since I was little}

\label{lesson:ZH-C269-congxiao}



\begin{quote}
Before the new one: say the Mandarin for the answer to a question, then the Mandarin for to explain.
\end{quote}

\subsection*{You'll want to know: \zh{从小}}
\textbf{\zh{从小}} — \emph{cóngxiǎo} — \textquotedblleft{}since childhood; since I was little\textquotedblright{}.

\begin{quote}
\textbf{\zh{我从小喜欢看书}} \emph{wǒ cóngxiǎo xǐhuan kànshū} — I've loved books since I was little.
\end{quote}

\begin{cousinweb}
\textbf{\zh{从小}} (\emph{cóngxiǎo}) is built from \zh{从} 'from' and \zh{小} 'small', both inherited from Middle Chinese, \emph{dzjowng} and \emph{sjewX} (Guangyun, in Baxter's transcription).
\end{cousinweb}

\begin{grammarlens}[title={What you've built}]
One more word for talking about the past.
\end{grammarlens}

\subsection*{Your turn}
\begin{itemize}
\raggedright
  \item \emph{Say it:} \emph{cóngxiǎo}
  \item \emph{Say it:} \emph{cóngxiǎo}, once more
  \item \emph{Say it:} \emph{wǒ cóngxiǎo xǐhuan kànshū} — I've loved books since I was little.
  \item \emph{Recall:} say the Mandarin for the answer to a question, then the Mandarin for to explain, then say \emph{cóngxiǎo} again
\end{itemize}

\begin{tcolorbox}[breakable,colback=teal!4,colframe=teal!35!black,title={Before you move on}]
What does \zh{从小} mean? (\textquotedblleft{}Since childhood; since I was little\textquotedblright{}.) Say it once more.
\end{tcolorbox}
\section[\texorpdfstring{\zh{先} (xiān)}{先 (xiān)}]{\zh{先} (xiān) — first; before anything else}

\label{lesson:ZH-C269-xian}



\begin{quote}
Before the new one: say the Mandarin for to explain, then the Mandarin for since childhood.
\end{quote}

\subsection*{You'll want to know: \zh{先}}
\textbf{\zh{先}} — \emph{xiān} — \textquotedblleft{}first; before anything else\textquotedblright{}.

\begin{quote}
\textbf{\zh{我们先吃饭}} \emph{wǒmen xiān chīfàn} — Let's eat first.
\end{quote}

The new character, in pieces: \zh{先} — six strokes: a short stroke falling to the left, a short bar, a vertical and a long bar; then \zh{儿} at the foot. The next lesson puts it in your hand.

\begin{cousinweb}
\textbf{\zh{先}} (\emph{xiān}) is inherited from Middle Chinese \emph{sen} (Guangyun, in Baxter's transcription), glossed 'before and after'.
\end{cousinweb}

\begin{grammarlens}[title={What you've built}]
One more word for talking about the past.
\end{grammarlens}

\subsection*{Your turn}
\begin{itemize}
\raggedright
  \item \emph{Say it:} \emph{xiān}
  \item \emph{Say it:} \emph{xiān}, once more
  \item \emph{Say it:} \emph{wǒmen xiān chīfàn} — Let's eat first.
  \item \emph{Recall:} say the Mandarin for to explain, then the Mandarin for since childhood, then say \emph{xiān} again
\end{itemize}

\begin{tcolorbox}[breakable,colback=teal!4,colframe=teal!35!black,title={Before you move on}]
What does \zh{先} mean? (\textquotedblleft{}First; before anything else\textquotedblright{}.) Say it once more.
\end{tcolorbox}
\section[\texorpdfstring{\zh{先} (xiān)}{先 (xiān)}]{\zh{先} — a cross and a long bar, over \zh{儿}}

\label{lesson:ZH-W269-xian-first}



\begin{quote}
The word first, then the character inside it.

\begin{itemize}
\raggedright
  \item \emph{Recall:} say \emph{xiān} — \textbf{R1}, one lesson back, before it decays
\end{itemize}

\begin{quote}
\zh{儿} \zh{生}
\end{quote}

\emph{Read it:} each one aloud, with nothing said first — \emph{ér}, \emph{shēng}
\end{quote}

\subsection*{Script — a cross and a long bar, over \zh{儿}}
\hlblockfigure{\detokenize{figures/ZH-W269-xian-first-filmstrip.pdf}}{How \zh{先} is written, stroke by stroke}

\begin{quote}
\zh{先}
\end{quote}

\textbf{Two pieces, top before bottom.}

\textbf{The top piece, four strokes.} A short stroke falling to the left. A short bar to the right of it. A vertical down through the bar. Then a long bar across, wider than everything above it.

\textbf{The bottom piece, \zh{儿}, two strokes,} as on its own: a stroke falling to the left, then a vertical that bends to the right along the foot and hooks up.

4 + 2 = six.

\subsection*{Your turn}
\begin{itemize}
\raggedright
  \item \emph{Write it:} \zh{先}, all six, with the model in front of you
  \item \emph{Cover:} write it again from memory, and count the pieces, not the strokes
\end{itemize}

\begin{tcolorbox}[breakable,colback=teal!4,colframe=teal!35!black,title={Before you move on}]
How many strokes in \zh{先}? (\textbf{Six}.) What sits at the foot? (\textbf{\zh{儿}}.) Which bar is longer? (\textbf{The second, across the middle}.)

Source: \href{https://raw.githubusercontent.com/chanind/hanzi-writer-data/68d10a4b21150cae5e1ebbd223eed289cf32d90c/data/\%E5\%85\%88.json}{Hanzi Writer Data \zh{先}, pinned stroke paths and medians}.
\end{tcolorbox}
\section[tóngnián]{tóngnián — childhood}

\label{lesson:ZH-C269-tongnian}



\begin{quote}
Before the new one: say the Mandarin for since childhood, then the Mandarin for before anything else.
\end{quote}

\subsection*{You'll want to know: tóngnián}
\emph{tóngnián} — \textquotedblleft{}childhood\textquotedblright{}.

\begin{quote}
\emph{wǒ de tóngnián hěn kuàilè} — I had a very happy childhood.
\end{quote}

Its characters come later; for now it is a word for the ear.

\begin{cousinweb}
\emph{tóngnián} is built from \emph{tóng} 'a child' and \zh{年} 'a year', both inherited from Middle Chinese, \emph{duwng} and \emph{nen} (Guangyun, in Baxter's transcription), which glosses the first as a young person who does not yet have a household.
\end{cousinweb}

\begin{grammarlens}[title={What you've built}]
One more word for talking about the past.
\end{grammarlens}

\subsection*{Your turn}
\begin{itemize}
\raggedright
  \item \emph{Say it:} \emph{tóngnián}
  \item \emph{Say it:} \emph{tóngnián}, once more
  \item \emph{Say it:} \emph{wǒ de tóngnián hěn kuàilè} — I had a very happy childhood.
  \item \emph{Recall:} say the Mandarin for since childhood, then the Mandarin for before anything else, then say \emph{tóngnián} again
\end{itemize}

\begin{tcolorbox}[breakable,colback=teal!4,colframe=teal!35!black,title={Before you move on}]
What does tóngnián mean? (\textquotedblleft{}Childhood\textquotedblright{}.) Say it once more.
\end{tcolorbox}
\section[\texorpdfstring{\zh{本来} (běnlái)}{本来 (běnlái)}]{\zh{本来} (běnlái) — originally; at first}

\label{lesson:ZH-C269-benlai}



\begin{quote}
Before the new one: say the Mandarin for before anything else, then the Mandarin for childhood.
\end{quote}

\subsection*{You'll want to know: \zh{本来}}
\textbf{\zh{本来}} — \emph{běnlái} — \textquotedblleft{}originally; at first\textquotedblright{}.

\begin{quote}
\textbf{\zh{我本来想先吃饭}} \emph{wǒ běnlái xiǎng xiān chīfàn} — I'd meant to eat first.
\end{quote}

\begin{cousinweb}
\textbf{\zh{本来}} (\emph{běnlái}) is built from \zh{本} 'the root' and \zh{来} 'to come', both inherited from Middle Chinese, \emph{pwonX} and \emph{loj} (Guangyun, in Baxter's transcription), which glosses the first 'the root' and 'the beginning'.
\end{cousinweb}

\begin{grammarlens}[title={What you've built}]
One more word for talking about the past.
\end{grammarlens}

\subsection*{Your turn}
\begin{itemize}
\raggedright
  \item \emph{Say it:} \emph{běnlái}
  \item \emph{Say it:} \emph{běnlái}, once more
  \item \emph{Say it:} \emph{wǒ běnlái xiǎng xiān chīfàn} — I'd meant to eat first.
  \item \emph{Recall:} say the Mandarin for before anything else, then the Mandarin for childhood, then say \emph{běnlái} again
\end{itemize}

\begin{tcolorbox}[breakable,colback=teal!4,colframe=teal!35!black,title={Before you move on}]
What does \zh{本来} mean? (\textquotedblleft{}Originally; at first\textquotedblright{}.) Say it once more.
\end{tcolorbox}
\section[wánjù]{wánjù — toy}

\label{lesson:ZH-C269-wanju}



\begin{quote}
Before the new one: say the Mandarin for childhood, then the Mandarin for originally.
\end{quote}

\subsection*{You'll want to know: wánjù}
\emph{wánjù} — \textquotedblleft{}toy\textquotedblright{}.

\begin{quote}
\emph{zhège wánjù hěn yǒu yìsi} — This toy is a lot of fun.
\end{quote}

Its characters come later; for now it is a word for the ear.

\begin{cousinweb}
\emph{wánjù} is built from \emph{wán} 'to play' and \emph{jù}, both inherited from Middle Chinese, \emph{ngwanH} and \emph{gjuH} (Guangyun, in Baxter's transcription), which glosses the first 'to toy with' and the second 'to provide'. Middle Chinese \emph{ngwanH} would regularly give wàn; the rising tone of wán is irregular.
\end{cousinweb}

\begin{grammarlens}[title={What you've built}]
One more word for talking about the past.
\end{grammarlens}

\subsection*{Your turn}
\begin{itemize}
\raggedright
  \item \emph{Say it:} \emph{wánjù}
  \item \emph{Say it:} \emph{wánjù}, once more
  \item \emph{Say it:} \emph{zhège wánjù hěn yǒu yìsi} — This toy is a lot of fun.
  \item \emph{Recall:} say the Mandarin for childhood, then the Mandarin for originally, then say \emph{wánjù} again
\end{itemize}

\begin{tcolorbox}[breakable,colback=teal!4,colframe=teal!35!black,title={Before you move on}]
What does wánjù mean? (\textquotedblleft{}Toy\textquotedblright{}.) Say it once more.
\end{tcolorbox}
